\section{Síntesis y ampliación}

Esta sesión concretó el test-time compute, que antes era «pensar un poco más», en dos sistemas ejecutables.

\begin{enumerate}
    \item \textbf{AlphaCode}: cubre el espacio de programas mediante muestreo masivo en paralelo, y comprime los candidatos usando casos de ejemplo, agrupamiento por comportamiento y pruebas ocultas;
    \item \textbf{AlphaCode 2}: mejora la eficiencia de muestreo y la calidad de las presentaciones con un foundation model más potente, múltiples solver family y un scorer aprendido;
    \item \textbf{Search-o1}: identifica vacíos de conocimiento dentro del reasoning, recupera y comprime documentos en varias rondas, y deja que la nueva evidencia entre al razonamiento posterior;
    \item \textbf{Search-R1}: trata «cuándo buscar, qué buscar y cuándo detenerse» como una política que se puede aprender mediante RL.
\end{enumerate}

Al recorrer tanto la búsqueda de código como la búsqueda de conocimiento, se puede destilar el mismo conjunto de principios de diseño de agent:

\begin{itemize}
    \item el generador debe ampliar la cobertura de candidatos, pero manteniendo diversidad semántica;
    \item el selector debe basarse en señales verificables relevantes para la tarea, no en la probabilidad de generación;
    \item el presupuesto de cómputo debe asignarse dinámicamente según la dificultad y la incertidumbre actual;
    \item las trayectorias largas deben conservar un estado intermedio estructurado que permita la verificación local y el retroceso;
    \item las herramientas externas introducen nuevas superficies de error, por lo que los resultados de la búsqueda, las pruebas y el propio reward model también deben verificarse.
\end{itemize}

La siguiente sesión se dirige hacia la agentic evaluation: cómo medir la confiabilidad, el valor económico y la brecha con el trabajo real cuando la tarea requiere de decenas de minutos a varias horas y la salida es un CAD, un informe, una tabla o una revisión de literatura.

\end{document}
